% page 10 du fac-simile (image p-009 du scan)
% bloc 157.5 x 233.3 mm ; marge gauche 8.0 mm
\pgc{-12.700mm}{0.000mm}{
\l{\cel{4}2}
\l{}
\l{\sou{abciso}. (\sou{geom.)}\cel{2}Un ek la du koordinati rekta-linea qui}
\l{\cel{3}uzesas por determinar la situeso di punto en spaco.-}
\l{\cel{3}DEFIRS.}
\l{}
\l{\sou{abdikar}. (\sou{trans}.) (\sou{pri rejo od altra suvereno)}. Cesar sua}
\l{\cel{3}regno per transmisar la autoritato imperala a sucedanto,}
\l{\cel{3}o naturala, o selektita.\cel{2}- DEFIRS.}
\l{}
\l{\sou{abdomino}.\cel{2}(\sou{anat}.)Korpo-kavajo qua kontenas la viceri digest-}
\l{\cel{3}tala. (Ta vorto ciencala esas plu preciza kam \textquotedbl{}ventro\textquotedbl{}).}
\l{\cel{3}- DEFIS.}
\l{}
\l{\sou{abduktar}.-\sou{trans}.) Movigar parto di la homo-korpo tale ke la}
\l{\cel{3}movajo eskartesas de la plano mediana di la korpo.-DEFIS.}
\l{}
\l{\sou{abelo}. (\sou{zool}.) Insekto himenoptera qua vivas esame e produk-}
\l{\cel{3}tas mielo e vaxo. - L.\cel{1}\sou{apis mellifica}.-DEFIS.}
\l{}
\l{\sou{aberaco}. I. (\sou{astron.})\cel{2}Deviaceso semblanta di la radii luma}
\l{\cel{3}qui venas de astro. - II. (\sou{fiziko}) Disperseso di la radii}
\l{\cel{3}luma qui trairas lenso. - DEFIRS.}
\l{}
\l{\sou{abieto}. (\sou{bot}.) Arboro rezinifera, di qua la folii nula-sezone}
\l{\cel{3}divenas flavatra e sika. - L.\cel{1}\sou{abies pectinata}.- EFIS.}
\l{}
\l{\sou{abiso}.-Profundegajo submara.}
\l{}
\l{\sou{abismo}. I. Profundegajo qua ne esas mezurebla pro la grand-}
\l{\cel{3}eso di lua dimensioni. - II. (\sou{metaf}.)(a) Profundegajo quan}
\l{\cel{3}spirito ne povas mezurar konceptale. (b) Profundegajo en}
\l{\cel{3}qua onu esas quaze nihiligata, su-egaranta. - EFIS.}
\l{}
\l{\sou{abjekta}. Qua esas ye la grado maxim infra di desnobleso, pri}
\l{\cel{3}la hierarkio sociala e pri lua karaktero personala.-DEFIS}
\l{}
\l{\sou{abjurar}. (\sou{trans.)}\cel{1}I. Renuncar, en su-obligo solena, la religio}
\l{\cel{3}quan onu profesis til-lore. -II. (\sou{metaf}.)Renuncar publike}
\l{\cel{3}to quon onu profesis kredar, amar, prizar.- DEFIS.}
\l{}
\l{\sou{ablacionar}. (\sou{trans}.)(\sou{kirurg.}) Deprenar de la korpo homala}
\l{\cel{3}(per operaco kirurgiala) parto morboza, malada.-DEFIS.}
\l{}
\l{\sou{ablaktar}. (\sou{trans.}) Separar infanteto de lua patrino o de la}
\l{\cel{3}salariato qua alaktis lu, kande oportas komencar donar}
\l{\cel{3}a lu alimente nutrivi altra kam lakto. - DEFIS.}
\l{}
\l{\sou{ablativo}. (\sou{gram.}) Kazo di la deklino-sistemo Greka, Latina,}
\l{\cel{3}e c., qua indikas ke substantivo esas la departo-punto,}
\l{\cel{3}o la instrumento, di ago. - DEFIRS.}
\l{}
\l{\sou{ablegato}. (\sou{religio katolika}). Delegito qua komisesis por por-}
\l{\cel{3}tar la bireto a prelato jus nominita \textquotedbl{}kardinalo\textquotedbl{}.-DEFI.}
}
